% !TeX root = ../../Book4.tex
% 纲要标题：### 8.3 马蹄引理
% 纲要标签：b4:sec:0703
% 纲要位置：计划纲要.md，第 3287 行。

\section{马蹄引理}
\label{b4:sec:0703}

短正合列的两端已有消解时，中项的消解仍须保存原扩张的不分裂信息。马蹄构造把这部分信息放进微分的非对角项，使三份消解逐次相容。


% 纲要内容（仅作写作计划，尚非教材正文）：
%
% * 短正合列与相容消解
% * 投射及内射版本
% * 消解上的长正合列
%

\subsection{短正合列与相容消解}
\label{b4:subsec:0703:01}

假设已有短正合列 \(0\to A\xrightarrow{i}B\xrightarrow{q}C\to0\)，又已分别选好 \(A,C\) 的投射消解。希望构造 \(B\) 的消解时，最自然的逐次对象是 \(P_n^A\oplus P_n^C\)。困难在于微分通常不能取两个微分的直和：那样得到的零次同调是 \(A\oplus C\)，未必是 \(B\)。原列的不分裂信息必须进入微分的非对角项。

\begin{lemma}\label{b4:lem:horseshoe-step}
设 \(\mathcal A\) 是交换范畴，\(0\to A\xrightarrow{i}B\xrightarrow{q}C\to0\) 正合，\(p_A:P_A\to A\)、\(p_C:P_C\to C\) 满，且 \(P_C\) 投射。选 \(t:P_C\to B\) 使 \(qt=p_C\)，令 \(p_B=(ip_A,t):P_A\oplus P_C\to B\)。则 \(p_B\) 满，且
\[
0\longrightarrow\ker p_A\longrightarrow\ker p_B
\longrightarrow\ker p_C\longrightarrow0
\]
正合。
\end{lemma}
\begin{proof}
在两行分别为 \(0\to P_A\to P_A\oplus P_C\to P_C\to0\) 和给定短正合列的图中，竖直态射取 \(p_A,p_B,p_C\)。图由 \(qt=p_C\)、\(qi=0\) 交换：
\[
\begin{tikzcd}[column sep=2.5em,row sep=large]
0\arrow[r]&P_A\arrow[r]\arrow[d,two heads,"p_A"']&P_A\oplus P_C\arrow[r]\arrow[d,"p_B"]&P_C\arrow[r]\arrow[d,two heads,"p_C"]&0\\
0\arrow[r]&A\arrow[r,"i"']&B\arrow[r,"q"']&C\arrow[r]&0
\end{tikzcd}
\]
由\cref{b4:thm:snake-lemma}得到正合列
\[
\begin{aligned}
0\longrightarrow\ker p_A\longrightarrow\ker p_B
&\longrightarrow\ker p_C\longrightarrow\operatorname{coker}p_A\\
&\longrightarrow\operatorname{coker}p_B
\longrightarrow\operatorname{coker}p_C\longrightarrow0.
\end{aligned}
\]
因 \(p_A,p_C\) 均满，两端余核为零。于是 \(\ker p_B\to\ker p_C\) 满，且 \(\operatorname{coker}p_B=0\)，这正是结论。
\end{proof}

\subsection{马蹄引理的投射与内射版本}
\label{b4:subsec:0703:02}

\begin{theorem}[马蹄引理]\label{b4:thm:horseshoe}
设 \(\mathcal A\) 是交换范畴，\(0\to A\to B\to C\to0\) 是短正合列。给定 \(A,C\) 的投射消解 \(P^A_\bullet,P^C_\bullet\)，存在 \(B\) 的投射消解 \(P^B_\bullet\)，使增广相容，且
\[
0\longrightarrow P^A_\bullet\longrightarrow P^B_\bullet
\longrightarrow P^C_\bullet\longrightarrow0
\]
逐次分裂正合，并可取 \(P_n^B=P_n^A\oplus P_n^C\)。对偶地，给定 \(A,C\) 的内射消解，可取 \(I_B^n=I_A^n\oplus I_C^n\)，得到相容且逐次分裂的短正合列 \(0\to I_A^\bullet\to I_B^\bullet\to I_C^\bullet\to0\)。此结果称为\term[matiyinli]{马蹄引理}[horseshoe lemma]。
\end{theorem}
\begin{proof}
这里的分裂逐次数成立，截面未必与微分交换。选定双积表示后，中间微分的非对角块正是需要另外构造的部分；只有存在复形截面时，才能用它选取双积表示，使这些非对角块同时为零。

先将\cref{b4:lem:horseshoe-step}用于 \(P_0^A\to A\)、\(P_0^C\to C\)，得到 \(P_0^B\to B\) 及其核列。记这三个核为 \(K_1^A,K_1^B,K_1^C\)。给定消解中的下一项分别满射到 \(K_1^A,K_1^C\)，于是再用引理得到 \(P_1^B\to K_1^B\)，其核又组成短正合列。归纳继续。把 \(P_n^B\to K_n^B\) 与 \(K_n^B\hookrightarrow P_{n-1}^B\) 复合，就是中间消解的微分；每一步定义保证像等于下一个核。

各次包含与投影通过核图构造，因而与微分交换。其对象列是双积的分裂列；\(P_n^A\oplus P_n^C\) 投射，因为两个分量的态射可以分别提升后相加。这样得到全部结论。内射版本在 \(\mathcal A^{\mathrm{op}}\) 应用已证投射版本，并反转箭头与链指标；有限双积在反范畴仍是双积，故逐次分裂性不变。
\end{proof}

选定逐次双积表示以后，相容消解可以排列如下：
\[
\begin{tikzcd}[column sep=large,row sep=large]
P_1^A\arrow[r]\arrow[d,"\mathrm{d}_1^A"']&P_1^B\arrow[r]\arrow[d,"\mathrm{d}_1^B"]&P_1^C\arrow[d,"\mathrm{d}_1^C"]\\
P_0^A\arrow[r]\arrow[d]&P_0^B\arrow[r]\arrow[d]&P_0^C\arrow[d]\\
A\arrow[r]&B\arrow[r]&C
\end{tikzcd}
\]
从 \(P_0^A,P_0^B,P_0^C\) 向下的箭头为增广；上面两排来自逐次分裂的消解短正合列。微分具体为
\[
\mathrm{d}_n^B=\begin{pmatrix}\mathrm{d}_n^A&h_n\\0&\mathrm{d}_n^C\end{pmatrix},
\qquad \mathrm{d}_{n-1}^Ah_n+h_{n-1}\mathrm{d}_n^C=0.
\]
第二个等式在 \(n\ge2\) 时成立，是 \((\mathrm{d}^B)^2=0\) 的非对角分量，显示原来的扩张如何影响消解。零次的相容关系则由增广给出。

\subsection{消解上的长正合列}
\label{b4:subsec:0703:03}

\begin{proposition}\label{b4:prop:horseshoe-functor-sequence}
设 \(\mathcal A,\mathcal B\) 为交换范畴，\(F:\mathcal A\to\mathcal B\) 为加性函子，\(0\to A\to B\to C\to0\) 为 \(\mathcal A\) 中的短正合列。给定两端 \(A,C\) 的投射消解，并按马蹄引理构造相容的消解短正合列 \(0\to P^A_\bullet\to P^B_\bullet\to P^C_\bullet\to0\)。此列经 \(F\) 逐次作用后仍是逐次分裂的复形短正合列，因此其同调对象有自然连接态射及长正合列。两端取内射消解及相应的马蹄构造时，同样成立。
\end{proposition}
\begin{proof}
分裂短正合列由态射等式 \(ri=1\)、\(ps=1\)、\(ir+sp=1\) 表示。加性函子保持这些等式，故不必假定 \(F\) 正合。再用\cref{b4:thm:cohomology-long-exact}，在链情形反转指标，即得到长正合列。

为证明自然性，给定短正合列间的交换图
\[
\begin{tikzcd}[column sep=large,row sep=large]
0\arrow[r]&A\arrow[r,"i"]\arrow[d,"a"']
&B\arrow[r,"q"]\arrow[d,"b"]
&C\arrow[r]\arrow[d,"c"]&0\\
0\arrow[r]&A'\arrow[r,"i'"']
&B'\arrow[r,"q'"']&C'\arrow[r]&0.
\end{tikzcd}
\]
源端的马蹄消解记为 \(P^A,P^B,P^C\)，目标端记为 \(Q^A,Q^B,Q^C\)。选定分别提升 \(a,c\) 的端点比较映射 \(u:P^A\to Q^A\)、\(v:P^C\to Q^C\)。用双积表示两套微分，分别记非对角项为 \(h_n\)、\(h'_n\)。设各端点增广为 \(\varepsilon_{P^A},\varepsilon_{P^C},\varepsilon_{Q^A},\varepsilon_{Q^C}\)，中间增广写成
\[
\begin{aligned}
\varepsilon_P&=(i\varepsilon_{P^A},t_P):
P_0^A\oplus P_0^C\longrightarrow B,\\
\varepsilon_Q&=(i'\varepsilon_{Q^A},t_Q):
Q_0^A\oplus Q_0^C\longrightarrow B',
\end{aligned}
\]
其中 \(qt_P=\varepsilon_{P^C}\)、\(q't_Q=\varepsilon_{Q^C}\)。我们构造中间比较映射
\[
w_n=\begin{pmatrix}u_n&y_n\\0&v_n\end{pmatrix}.
\]
零次需要满足 \(\varepsilon_Qw_0=b\varepsilon_P\)。第一列已由 \(u_0\) 的增广条件满足；第二列要求
\[
i'\varepsilon_{Q^A}y_0=bt_P-t_Qv_0:
P_0^C\longrightarrow B'.
\]
右端与 \(q'\) 的复合为
\[
q'(bt_P-t_Qv_0)
=c\varepsilon_{P^C}-\varepsilon_{Q^C}v_0=0.
\]
因 \(i'\) 是 \(q'\) 的核，右端唯一分解为 \(i'z_0\)，其中 \(z_0:P_0^C\to A'\)。由 \(P_0^C\) 投射，沿满态射 \(\varepsilon_{Q^A}:Q_0^A\to A'\) 提升 \(z_0\)，得到 \(y_0:P_0^C\to Q_0^A\)。这就满足了中间的增广条件。

已构造到次数 \(n-1\) 后，次数 \(n\) 的链映射条件只剩非对角块
\[
\mathrm{d}_n^{Q^A}y_n=r_n,
\qquad r_n=u_{n-1}h_n+y_{n-1}\mathrm{d}_n^{P^C}-h'_nv_n.
\]
这里 \(r_n:P_n^C\to Q_{n-1}^A\)，其余三块由 \(u,v\) 是链映射及零块自动满足。当 \(n=1\) 时，源、目标增广与微分的相容关系分别给出
\[
i\varepsilon_{P^A}h_1+t_P\mathrm{d}_1^{P^C}=0,
\qquad
i'\varepsilon_{Q^A}h'_1+t_Q\mathrm{d}_1^{Q^C}=0.
\]
结合已构造的 \(y_0\)，可得
\[
\begin{aligned}
i'\varepsilon_{Q^A}r_1
&=-bt_P\mathrm{d}_1^{P^C}
 +(bt_P-t_Qv_0)\mathrm{d}_1^{P^C}
 +t_Q\mathrm{d}_1^{Q^C}v_1\\
&=0,
\end{aligned}
\]
最后一步使用 \(v_0\mathrm{d}_1^{P^C}=\mathrm{d}_1^{Q^C}v_1\)。因 \(i'\) 单，\(\varepsilon_{Q^A}r_1=0\)。当 \(n\ge2\) 时，利用两套微分的非对角关系和上一次数的链映射等式，得到
\[
\begin{aligned}
\mathrm{d}_{n-1}^{Q^A}r_n
={}&-u_{n-2}h_{n-1}\mathrm{d}_n^{P^C}\\
&+\bigl(u_{n-2}h_{n-1}
 +y_{n-2}\mathrm{d}_{n-1}^{P^C}
 -h'_{n-1}v_{n-1}\bigr)\mathrm{d}_n^{P^C}\\
&+h'_{n-1}\mathrm{d}_n^{Q^C}v_n
=0.
\end{aligned}
\]
其中含 \(y_{n-2}\) 的项由 \((\mathrm{d}^{P^C})^2=0\) 消失，其余项由 \(v\) 的链映射等式抵消。因此 \(r_n\) 分解到目标核；目标消解正合，使 \(Q_n^A\) 满射到该核，再由 \(P_n^C\) 投射取提升 \(y_n\)。归纳给出消解短正合列之间的严格交换图。

对图施加 \(F\)，再用复形长正合列的自然性，即有
\[
H_{n-1}(F(u))\,\delta_P
=\delta_Q\,H_n(F(v)).
\]
若固定这两套消解而改变比较映射或中间提升，则所得三个比较链映射仍分别提升 \(a,b,c\)。由\cref{b4:prop:comparison-homotopy-unique}，它们分别链同伦；加性函子保持这些同伦，故诱导的同调态射不变。

若改变端点消解或马蹄消解，则把上述构造用于原短正合列的恒等图。三个比较映射都提升恒等，由\cref{b4:cor:resolutions-homotopy-equivalent}在施加 \(F\) 后诱导同调同构。比较映射的同伦唯一性保证这些同构不依赖提升，并满足传递律；刚才的自然性等式保证它们也与连接态射交换。因此，不同消解所得的长正合列在这些典范比较同构下相容，而不是在字面上具有相同的同调对象。内射情形在反范畴作同一构造，转回后得到上同调长正合列及其相容性。
\end{proof}

连接态射有可操作的含义：把商复形中的圈提升到中间复形，对提升取微分，再把所得像读成左端复形中的圈。马蹄引理保证即使函子不正合，仍有逐次正合的复形列来执行这一过程；下一章的 Ext 与 Tor 正是由此获得正合列。
